\subsection{עקרונות ושיטות בעיבוד הנתונים}

לפני הזנת הנתונים למודלים, כל התמונות עוברות סדרה של פעולות עיבוד
מקדים (\textenglish{Preprocessing}) ואוגמנטציה
(\textenglish{Augmentation}). מטרת שלב זה היא לשפר את יכולת ההכללה
(\textenglish{Generalization}) של המודל ולהבטיח יציבות במהלך האימון.

\begin{itemize}
\item \textbf{נרמול (\textenglish{Normalization}):} תמונות הקלט
  מורכבות מפיקסלים בטווח [0, 255]. כחלק מהעיבוד המוקדם, כל ערכי
  הפיקסלים מחולקים ב־255 כדי להעביר אותם לטווח [0, 1]. פעולה זו קריטית
  למניעת תופעת הגרדיאנטים המתפוצצים (\textenglish{Exploding
    Gradients}), שבה ערכי קלט גבוהים גורמים לעדכוני משקולות קיצוניים
  המונעים מה־\textenglish{Optimizer} להתכנס למינימום בצורה
  יציבה. הנרמול מאפשר למידה מהירה ועקבית יותר.
  
\item \textbf{אוגמנטציה (\textenglish{Data Augmentation}):} כדי
  להתמודד עם כמות נתונים מוגבלת ומגוון תנאי צילום אפשריים (תאורה
  משתנה, זוויות שונות, טשטוש תנועה), השתמשתי בטכניקות אוגמנטציה. פעולה
  זו "מייצרת" דגימות חדשות מתוך הקיימות על ידי שינויים גאומטריים
  ואופטיים, מה שמאלץ את הרשת ללמוד מאפיינים עמידים
  (\textenglish{Robust Features}) ולא להסתמך על רעשי רקע.
\end{itemize}

\subsection{מודל זיהוי הרובוטים}

השתמשתי
ב\href{https://universe.roboflow.com/main-wcgiu/robot-detection-xru6m/16}{מערך
  נתונים} קיים מ־Roboflow, עם מעל 1000 תמונות משידורים רשמיים של
\textenglish{FIRST}. כל רובוט מסומן ב־\textenglish{Bounding Box} שלו
ומסווג לפי הברית שלו: כחול או אדום.

\begin{itemize}
  \item \textbf{אוגמנטציה:} על הנתונים יש אוגמנטציה שמכפילה את מספר תמונות האימון
  פי 3 כדי להגדיל את כמות הנתונים ולהוסיף גיוון כהכנה לתנאי צילום שונים:
  \begin{english}
    \begin{itemize}
    \item \textbf{Shear:} \qty{+-10}{\degree} Horizontal, \qty{+-10}{\degree} Vertical
    \item \textbf{Brightness:} \qty{+-20}{\percent}
    \item \textbf{Noise:} $\leq$ \qty{0.1}{\percent}
    \end{itemize}
  \end{english}

\item \textbf{עיבוד מקדים (\textenglish{Preprocessing}):} התמונה נמתחת
  לגודל של \textenglish{640x640}. התמונה נשארת צבעונית כדי לאפשר את
  סיווג הברית, ועוברת נרמול לטווח [0, 1] על ידי \textenglish{YOLO}.
  
\item \textbf{פיצול הנתונים:} הנתונים פוצלו למערכי אימון, ולידציה
  ובדיקה בחלוקה של
  \textenglish{\qty{75}{\percent}--\qty{15}{\percent}--\qty{10}{\percent}}
  בהתאמה.
\end{itemize}

\subsection{מודל זיהוי הספרות}

השתמשתי
ב\href{https://app.roboflow.com/new-workspace-pyain/nums-llmwb/13}{מערך
  נתונים} עם כ-500 תמונות של רובוטים שתייגתי ידנית בעזרת
\textenglish{Roboflow}. כל ספרה מסומנת ב־\textenglish{Bounding Box}
שלה ומסווגת. התמונות המוקדמות יותר נאספו ונחתכו ידנית ממקורות שונים
כולל סרטונים באיכות גבוהה יותר, בעוד שהמאוחרות יותר נחתכו אוטומטית על
ידי מודל זיהוי הרובוטים מסרטונים רשמיים כדי לדמות את הקלט האמיתי
שהמודל יקבל.

תיוג הנתונים במערך זה לקח כשעתיים.

\begin{itemize}

\item \textbf{אוגמנטציה:} על הנתונים יש אוגמנטציה שמכפילה את מספר תמונות האימון
  פי 3, כדי להגדיל את כמות הנתונים ולהוסיף גיוון כהכנה לתנאי צילום
  שונים, לרובוטים שמופיעים בזוויות שונות ולחיתוכים לא מדויקים של הרובוט:
  \begin{english}
    \begin{itemize}
    \item \textbf{Crop:} $\leq$ \qty{10}{\percent}
    \item \textbf{Rotation:} \qty{+-5}{\degree}
    \item \textbf{Shear:} \qty{+-10}{\degree} Horizontal,
      \qty{+-5}{\degree} Vertical
    \item \textbf{Brightness:} \qty{+-15}{\percent}
    \item \textbf{Exposure:} \qty{+-10}{\percent}
    \end{itemize}
  \end{english}

\item \textbf{עיבוד מקדים (\textenglish{Preprocessing}):} התמונה
  מוגדלת לגודל של \textenglish{128x128} באמצעות
  \textenglish{Letterboxing} (הוספת שוליים שחורים לשמירה על יחס
  גובה־רוחב). התמונה נשארת צבעונית כדי לאפשר זיהוי טוב ופשוט יותר של הבאמפר
  (שעליו הספרות) ועוברת נרמול לטווח [0, 1] על ידי \textenglish{YOLO}.

\item \textbf{פיצול הנתונים:} בגלל הכמות הקטנה יחסית, הנתונים פוצלו
  למערכי אימון, ולידציה ובדיקה בחלוקה של
  \textenglish{\qty{85}{\percent}--\qty{10}{\percent}--\qty{5}{\percent}}
  בהתאמה.
\end{itemize}

\subsection{דוגמאות מתוך מערך הנתונים}

כדי להמחיש את איכות התיוג ואת השונות בנתונים, צורפו דוגמאות ויזואליות
למערכי הנתונים (איורים \ref{fig:robots_sample} ו־\ref{fig:digits_sample}) בנספח
\ref{app:data_samples}.

\subsection{התפלגות מחלקות}

מערך הנתונים כולל איזון בין רובוטים מהברית הכחולה לאדומה, ודגימות של
כל הספרות (0--9) במגוון זוויות צילום ותנאי תאורה, כדי להבטיח הכללה
(\textenglish{Generalization}) טובה של המודל.

להלן התפלגויות הדגימות בשני מערכי הנתונים. ניתן לראות כי המערך מאוזן
יחסית, מה שמסייע במניעת הטיות (Bias) של המודל.

\begin{figure}[H]
  \centering
  \begin{minipage}{0.4\textwidth}
    \centering
    \begin{tikzpicture}
      \begin{axis}[
        width=0.8\textwidth,
        scale only axis,
        ybar,
        xlabel={\texthebrew{מחלקה}},
        ylabel={\texthebrew{כמות}},
        symbolic x coords={כחול, אדום},
        xtick=data,
        ymin=0,
        nodes near coords,
        ]
        \addplot coordinates {(כחול, 3701) (אדום, 3776)};
      \end{axis}
    \end{tikzpicture}
    \caption{התפלגות המחלקות של מודל הרובוטים.}
  \end{minipage}
  \hfill
  \begin{minipage}{0.4\textwidth}
    \centering
    \begin{tikzpicture}
      \begin{axis}[
        width=0.8\textwidth,
        scale only axis,
        ybar,
        xlabel={\texthebrew{מחלקה}},
        ylabel={\texthebrew{כמות}},
        symbolic x coords={0, 1, 2, 3, 4, 5, 6, 7, 8, 9},
        xtick=data,
        ymin=0,
        nodes near coords,
        ]
        \addplot coordinates {(0, 178) (1, 293) (2, 140) (3, 175) (4, 221) (5, 179) (6, 183) (7, 124) (8, 130) (9, 162)};
      \end{axis}
    \end{tikzpicture}
    \caption{התפלגות המחלקות של מודל הספרות.}
  \end{minipage}
\end{figure}
